% =========================================================
% Proof: Chain Rule (via Carathéodory's Theorem)
% Source: volume-iii/analysis/differentiation/notes/notes-caratheodory.tex
% Mode: dual
% =========================================================

\subsection*{Chain Rule}
\label{prf:chain-rule}

\begin{remark*}[Return]
\hyperref[thm:chain-rule]{$\leftarrow$ Back to Theorem (Chain Rule) in Notes}
\end{remark*}

% ---------------------------------------------------------
% Proof I — Learning Version
% ---------------------------------------------------------
\subsubsection*{Proof I \textemdash\ Learning Version (Carathéodory, unnamed $\Phi$)}
\label{prf:chain-rule-naive}

\begin{proof}
Let $J, I \subseteq \mathbb{R}$ be open intervals, let $f : J \to \mathbb{R}$
and $g : I \to \mathbb{R}$ be functions with $f(J) \subseteq I$, and let
$c \in J$. Assume $f$ is differentiable at $c$ and $g$ is differentiable at
$d := f(c)$. The claim is that $g \circ f$ is differentiable at $c$ and
$(g \circ f)'(c) = g'(f(c)) \cdot f'(c)$.

\ProofStep{1}{Linearize $f$ at $c$.}
Since $f$ is differentiable at $c$, by \ByThm{Carathéodory's Theorem} there
exists $\varphi : J \to \mathbb{R}$, continuous at $c$, such that
\[
  f(x) - f(c) = \varphi(x)(x - c)
  \quad \text{for all } x \in J, \qquad \varphi(c) = f'(c).
\]

\ProofStep{2}{Linearize $g$ at $f(c)$.}
Let $d = f(c)$. Since $g$ is differentiable at $d$, by
\ByThm{Carathéodory's Theorem} there exists $\psi : I \to \mathbb{R}$,
continuous at $d$, such that
\[
  g(y) - g(d) = \psi(y)(y - d)
  \quad \text{for all } y \in I, \qquad \psi(d) = g'(d).
\]

\ProofStep{3}{Compose the identities.}
Substitute $y = f(x)$ into the identity for $g$:
\[
  g(f(x)) - g(f(c)) = \psi(f(x))\bigl[f(x) - f(c)\bigr].
\]
Now substitute the identity for $f(x) - f(c)$ from Step 1:
\[
  g(f(x)) - g(f(c)) = \psi(f(x)) \cdot \varphi(x) \cdot (x - c).
\]

\ProofStep{4}{Take the limit.}
Since $f$ is continuous at $c$ (by \ByThm{Differentiability Implies
Continuity}) and $\psi$ is continuous at $f(c)$, the composition
$\psi \circ f$ is continuous at $c$ by \ByThm{Continuity of Composite
Functions}. Since $\varphi$ is continuous at $c$ from Step 1, the product
$\psi(f(x)) \cdot \varphi(x)$ is continuous at $c$ by \ByThm{Product Rule
for Continuous Functions}. Therefore:
\[
  \lim_{x \to c} \psi(f(x)) \cdot \varphi(x)
  = \psi(f(c)) \cdot \varphi(c)
  = g'(f(c)) \cdot f'(c).
\]

\ProofStep{5}{Conclude.}
For $x \neq c$, divide Step 3's identity by $(x - c)$:
\[
  \frac{(g \circ f)(x) - (g \circ f)(c)}{x - c}
  = \psi(f(x)) \cdot \varphi(x).
\]
Taking the limit as $x \to c$ gives $(g \circ f)'(c) = g'(f(c)) \cdot f'(c)$.
\end{proof}

% ---------------------------------------------------------
% Proof II — Professional Standard
% ---------------------------------------------------------
\subsubsection*{Proof II \textemdash\ Professional Standard (Carathéodory, named $\Phi$)}
\label{prf:chain-rule-direct}

\begin{proof}
Let $J, I \subseteq \mathbb{R}$ be open intervals, let $f : J \to \mathbb{R}$
and $g : I \to \mathbb{R}$ be functions with $f(J) \subseteq I$, and let
$c \in J$. Assume $f$ is differentiable at $c$ and $g$ is differentiable at
$d := f(c)$. The claim is that $g \circ f$ is differentiable at $c$ and
$(g \circ f)'(c) = g'(f(c)) \cdot f'(c)$.

\ProofStep{1}{Apply Carathéodory's theorem to $f$ at $c$.}
By \ByThm{Carathéodory's Theorem}, there exists $\varphi : J \to \mathbb{R}$,
continuous at $c$, such that
\[
  f(x) - f(c) = \varphi(x)(x - c)
  \quad \text{for all } x \in J, \qquad \varphi(c) = f'(c).
\]

\ProofStep{2}{Apply Carathéodory's theorem to $g$ at $d$.}
By \ByThm{Carathéodory's Theorem}, there exists $\psi : I \to \mathbb{R}$,
continuous at $d = f(c)$, such that
\[
  g(y) - g(d) = \psi(y)(y - d)
  \quad \text{for all } y \in I, \qquad \psi(d) = g'(d).
\]

\ProofStep{3}{Compose the two factorizations.}
Substitute $y = f(x)$ into the identity for $g$, using $d = f(c)$:
\[
  g(f(x)) - g(f(c)) = \psi(f(x))\bigl[f(x) - f(c)\bigr].
\]
Substitute the identity for $f(x) - f(c)$ from Step 1:
\[
  g(f(x)) - g(f(c)) = \psi(f(x)) \cdot \varphi(x) \cdot (x - c).
\]

\ProofStep{4}{Read off the Carathéodory function for $g \circ f$.}
Define $\Phi : J \to \mathbb{R}$ by $\Phi(x) := \psi(f(x)) \cdot \varphi(x)$.
The identity from Step 3 reads
\[
  (g \circ f)(x) - (g \circ f)(c) = \Phi(x)(x - c)
  \quad \text{for all } x \in J.
\]
By \ByThm{Carathéodory's Theorem}, $g \circ f$ is differentiable at $c$ if
and only if $\Phi$ is continuous at $c$.

\ProofStep{5}{Verify continuity of $\Phi$ at $c$.}
Differentiability of $f$ at $c$ implies continuity of $f$ at $c$ by
\ByThm{Differentiability Implies Continuity}. Since $\psi$ is continuous at
$d = f(c)$ and $f$ is continuous at $c$, the composition $\psi \circ f$ is
continuous at $c$ by \ByThm{Continuity of Composite Functions}. Since
$\varphi$ is continuous at $c$ by Step 1, the product $\Phi = (\psi \circ f)
\cdot \varphi$ is continuous at $c$ by \ByThm{Product Rule for Continuous
Functions}.

\ProofStep{6}{Conclude and compute the derivative.}
By \ByThm{Carathéodory's Theorem}, $g \circ f$ is differentiable at $c$ and
$(g \circ f)'(c) = \Phi(c)$. Evaluating:
\[
  (g \circ f)'(c)
  = \Phi(c)
  = \psi(f(c)) \cdot \varphi(c)
  = g'(f(c)) \cdot f'(c). \qedhere
\]
\end{proof}

% ---------------------------------------------------------
% Shared remark blocks
% ---------------------------------------------------------
\begin{remark*}[Proof shape]
Both $f$ and $g$ are linearized via Carathéodory's theorem, yielding
Carathéodory functions $\varphi$ and $\psi$. Composing the two factorizations
produces a single Carathéodory factorization $\Phi = (\psi \circ f) \cdot
\varphi$ for $g \circ f$. Continuity of $\Phi$ at $c$ — established by
combining continuity of $\psi \circ f$ (via differentiability of $f$ and
continuity of $\psi$) with continuity of $\varphi$ — is then exactly the
Carathéodory criterion for differentiability of $g \circ f$. The derivative
formula falls out by evaluating $\Phi$ at $c$.

\FlashProof{1. Carathéodory on $f$ at $c$: get $\varphi$ continuous,
  $f(x)-f(c)=\varphi(x)(x-c)$, $\varphi(c)=f'(c)$.
  2. Carathéodory on $g$ at $f(c)$: get $\psi$ continuous,
  $g(y)-g(f(c))=\psi(y)(y-f(c))$, $\psi(f(c))=g'(f(c))$.
  3. Compose: $g(f(x))-g(f(c)) = \psi(f(x))\cdot\varphi(x)\cdot(x-c)$.
  4. Set $\Phi = (\psi\circ f)\cdot\varphi$; this is the Carathéodory
  function for $g\circ f$.
  5. $\Phi$ is continuous at $c$: diff implies cont for $f$, compose
  with $\psi$, multiply by $\varphi$.
  6. $(g\circ f)'(c) = \Phi(c) = g'(f(c))\cdot f'(c)$.}
\end{remark*}

\begin{remark*}[Self-assessment of Proof I]
Proof I is correct. The Carathéodory strategy is the right one, both
factorizations are correctly applied, the continuity chain is assembled, and
the limit is correctly computed. The following notes identify the structural
differences that separate a good proof from a professional one.

\textbf{The core gap: not naming the Carathéodory function for $g \circ f$.}
In Proof I, Steps 3 and 4 are merged into a single step that factors the
increment and then immediately passes to the limit. The product
$\psi(f(x)) \cdot \varphi(x)$ is treated as an expression to evaluate rather
than as an object to name and study. Proof II names this product $\Phi$ and
identifies it explicitly as the Carathéodory function for $g \circ f$. This
is not cosmetic. Naming $\Phi$ makes the proof's logical structure visible:
the question is whether $\Phi$ is continuous at $c$, which is exactly the
Carathéodory criterion applied to $g \circ f$. Without the name, the
continuity argument in Step 5 appears as a free-standing calculation; with
the name, it is clearly answering a specific question posed in Step 4.

\textbf{Step ordering: limit before difference quotient.}
Proof I takes the limit in Step 4 and then divides by $(x-c)$ in Step 5 to
exhibit the difference quotient. Proof II reverses this: Step 4 identifies
$\Phi$ as the Carathéodory function (which implicitly gives the difference
quotient), Step 5 establishes continuity of $\Phi$, and Step 6 applies the
theorem to conclude differentiability and read off the derivative. Proof II's
order is more logical: first identify what needs to be shown (continuity of
$\Phi$), then show it, then conclude. Proof I's order is more computational:
arrive at the answer, then justify it. The logical order is preferable because
it makes the proof's goal at each step clear before the work begins.

\textbf{Missed reuse of Carathéodory's theorem.}
Proof I applies Carathéodory in Steps 1 and 2 to linearize $f$ and $g$, but
never explicitly invokes it in the reverse direction to conclude
differentiability of $g \circ f$. The conclusion in Step 5 passes directly
from "the limit equals $g'(f(c)) \cdot f'(c)$" to "$(g \circ f)'(c) =
g'(f(c)) \cdot f'(c)$" by the definition of the derivative. This is correct
but misses an opportunity: Carathéodory's theorem was the tool that set up
the proof, and it is also the tool that closes it. Proof II explicitly
invokes it a third time in Step 6: "By Carathéodory's Theorem, $g \circ f$
is differentiable at $c$ and $(g \circ f)'(c) = \Phi(c)$." Symmetry of
structure — open with Carathéodory, close with Carathéodory — is a mark of
a well-organised proof.

\textbf{What Proof I does well.}
The Carathéodory strategy is independently chosen and correctly executed.
The substitution of $y = f(x)$ in Step 3 is done correctly. The continuity
chain — differentiability implies continuity of $f$, composition with $\psi$,
product with $\varphi$ — is assembled in the right order and with the right
theorem citations. The final evaluation $\psi(f(c)) \cdot \varphi(c) =
g'(f(c)) \cdot f'(c)$ is correct. All of this carries forward.
\end{remark*}

\begin{remark*}[Commentary]
\textbf{(a) Why use Carathéodory's theorem rather than the difference
quotient directly?}
The naive approach to the Chain Rule is to write the difference quotient of
$g \circ f$ as
\[
  \frac{g(f(x))-g(f(c))}{x-c}
  = \frac{g(f(x))-g(f(c))}{f(x)-f(c)} \cdot \frac{f(x)-f(c)}{x-c}
\]
and take limits separately. This fails when $f(x) = f(c)$ for $x \neq c$,
because the first factor becomes $0/0$. Carathéodory's theorem avoids
division entirely: the factorizations replace the quotients with continuous
functions, so the product $\psi(f(x)) \cdot \varphi(x)$ is well-defined
even when $f(x) = f(c)$.

\textbf{(b) Why apply Carathéodory to $g$ at $f(c)$, not at $c$?}
The function $g$ is differentiable at $d = f(c)$, not necessarily at $c$.
The Carathéodory function $\psi$ is therefore built around the point $d$:
it factorizes the increment $g(y) - g(f(c))$ by $(y - f(c))$. When $y = f(x)$,
this becomes $g(f(x)) - g(f(c))$ factorized by $f(x) - f(c)$, which is then
further factorized by $(x - c)$ via $\varphi$.

\textbf{(c) Why name $\Phi$?}
Naming $\Phi(x) := \psi(f(x)) \cdot \varphi(x)$ converts the factorization
\[
  (g \circ f)(x) - (g \circ f)(c) = \Phi(x)(x-c)
\]
into exactly the form that Carathéodory's theorem recognises. With $\Phi$
named, the entire remainder of the proof has a single question to answer: is
$\Phi$ continuous at $c$? Without the name, the continuity argument and the
conclusion are tangled together, making it harder to see what is being shown
and why.

\textbf{(d) Why does the continuity of $\Phi$ require three theorems?}
$\Phi = (\psi \circ f) \cdot \varphi$ is a product of two functions. Each
factor must be shown continuous at $c$ separately before their product can be
declared continuous. The factor $\psi \circ f$ is a composition: $\psi$ is
continuous at $f(c)$ (from Carathéodory applied to $g$), and $f$ is
continuous at $c$ (from differentiability of $f$), so their composition is
continuous at $c$ by the Continuity of Composite Functions theorem. The
factor $\varphi$ is continuous at $c$ directly from Carathéodory applied to
$f$. The Product Rule for Continuous Functions then combines the two. Each
theorem is doing a specific and necessary job.

\textbf{(e) Why does the derivative formula fall out at the end?}
Once $\Phi$ is shown to be continuous at $c$, Carathéodory's theorem gives
$(g \circ f)'(c) = \Phi(c)$. Evaluating:
$\Phi(c) = \psi(f(c)) \cdot \varphi(c) = g'(f(c)) \cdot f'(c)$.
The values $\psi(f(c)) = g'(f(c))$ and $\varphi(c) = f'(c)$ were recorded
at the end of Steps 2 and 1 respectively — they are the values of the
Carathéodory functions at the linearization points. The formula is not
derived; it is read off from what was already established.
\end{remark*}

\begin{remark*}[Dependencies]
\begin{itemize}
  \item \hyperref[thm:caratheodory]{Carathéodory's Theorem}: Applied to $f$
    at $c$ (Step 1), to $g$ at $f(c)$ (Step 2), and invoked in reverse to
    conclude differentiability of $g \circ f$ and read off the derivative
    (Step 6, Proof II).
  \item \hyperref[thm:diff-implies-cont]{Differentiability Implies
    Continuity}: Supplies continuity of $f$ at $c$ in Step 5, enabling the
    composition argument for $\psi \circ f$.
  \item \hyperref[thm:continuity-composition]{Continuity of Composite
    Functions}: Concludes $\psi \circ f$ is continuous at $c$ in Step 5.
  \item \hyperref[thm:continuity-product]{Product Rule for Continuous
    Functions}: Concludes $\Phi = (\psi \circ f) \cdot \varphi$ is
    continuous at $c$ in Step 5.
\end{itemize}

\FlashUses{thm:caratheodory, thm:diff-implies-cont, thm:continuity-composition, thm:continuity-product}
\end{remark*}